% Compiler avec XeLaTeX ou LuaLaTeX
\documentclass[11pt,a4paper]{article}

\usepackage{amsmath,amssymb}
\usepackage{fontspec}
\usepackage{unicode-math}
\usepackage{polyglossia}
\setmainlanguage{french}
\usepackage{xcolor}
\usepackage{fancyhdr}
\usepackage[margin=2cm,headheight=16pt,headsep=22pt]{geometry}
\usepackage{enumitem}

% Police principale et police mathématique : Libertinus
\IfFontExistsTF{Libertinus Serif}
  {\setmainfont{Libertinus Serif}}
  {\IfFontExistsTF{Latin Modern Roman}{\setmainfont{Latin Modern Roman}}{\setmainfont{lmroman10-regular.otf}}}

\IfFontExistsTF{Libertinus Math}
  {\setmathfont{Libertinus Math}}
  {\IfFontExistsTF{Latin Modern Math}{\setmathfont{Latin Modern Math}}{\setmathfont{latinmodern-math.otf}}}

\linespread{1.07}
\setlength{\parindent}{0pt}
\AtBeginDocument{\setlength{\parindent}{0pt}}
\setlength{\parskip}{0.35em}

% En-tête
\pagestyle{fancy}
\fancyhf{}
\lhead{Professeur de Mathématiques}
\rhead{Classe de Seconde}
% Pagination : « Page X / Y » centrée en pied de page
\usepackage{lastpage}
\cfoot{Page \thepage\ / \pageref{LastPage}}
\renewcommand{\headrulewidth}{0.5pt}
\renewcommand{\headrule}{\vskip 5pt\hrule width\headwidth height\headrulewidth \vskip-\headrulewidth}

% Configuration des listes
\setlist[enumerate,1]{itemsep=0.45em, topsep=0.3em, label=\textbf{\arabic*.}}
\setlist[enumerate,2]{itemsep=0.25em, topsep=0.2em, label=\textbf{\alph*.}}

% Commande exercice
\newcommand{\exo}[2]{%
  \vspace{1.5em}%
  \noindent{\large\textbf{Exercice #1. -- #2}}%
  \par\vspace{0.35em}%
  \hrule height 0.5pt%
  \vspace{0.7em}%
}

\newcommand{\partie}[1]{%
  \vspace{0.9em}%
  \noindent\textbf{#1}%
  \par\vspace{0.3em}%
}

\begin{document}

\begin{center}
  {\LARGE \textbf{Exercice d'entraînement sur les fractions}}\\[0.5em]
  {\large Classe de Seconde}
\end{center}

% ---------------------- EXERCICE 1 ----------------------
\exo{1}{Calculs et simplifications sur les fractions}

Cet exercice propose d'étudier des enchaînements d'opérations sur les fractions algébriques et numériques, en insistant sur les règles de réduction au même dénominateur, de factorisation et de simplification.

\partie{Partie A -- Calculs numériques}

Calculer les expressions suivantes et donner le résultat sous forme d'une fraction irréductible :
\begin{enumerate}
  \item $A = \frac{3}{4} + \frac{5}{6} \times \frac{2}{15}$
  \item $B = \left( \frac{7}{3} - \frac{1}{2} \right) \div \frac{5}{6}$
  \item $C = \frac{\frac{2}{3} + \frac{3}{4}}{\frac{5}{6} - \frac{1}{2}}$
\end{enumerate}

\partie{Partie B -- Simplification d'une expression littérale}

On considère l'expression $D$ définie pour tout réel $x \neq 0$ et $x \neq 2$ par :
$$D(x) = \frac{x^2 - 4}{x^2} \times \frac{x}{x - 2}$$

\begin{enumerate}
  \item Factoriser le numérateur $x^2 - 4$.
  \item En déduire que l'expression $D(x)$ peut s'écrire sous la forme simplifiée :
  $$D(x) = \frac{x + 2}{x}$$
  \item Calculer la valeur exacte de $D\left(\frac{1}{3}\right)$ sous forme de fraction irréductible.
\end{enumerate}

\vspace{1.5em}
\hrule height 0.5pt
\vspace{1em}
\noindent{\large\textbf{Corrigé détaillé}}

\partie{Correction de la Partie A}
\begin{enumerate}
  \item \textbf{Calcul de $A$ :}
  On effectue d'abord la multiplication en priorisant les opérations :
  $$\frac{5}{6} \times \frac{2}{15} = \frac{5 \times 2}{6 \times 15} = \frac{10}{90} = \frac{1}{9}$$
  Ensuite, on additionne avec $\frac{3}{4}$ en réduisant au même dénominateur ($36$) :
  $$A = \frac{3}{4} + \frac{1}{9} = \frac{3 \times 9}{4 \times 9} + \frac{1 \times 4}{9 \times 4} = \frac{27}{36} + \frac{4}{36} = \frac{31}{36}$$

  \item \textbf{Calcul de $B$ :}
  On calcule d'abord la différence entre parenthèses en réduisant au même dénominateur ($6$) :
  $$\frac{7}{3} - \frac{1}{2} = \frac{14}{6} - \frac{3}{6} = \frac{11}{6}$$
  Puis on divise par $\frac{5}{6}$ en multipliant par son inverse :
  $$B = \frac{11}{6} \div \frac{5}{6} = \frac{11}{6} \times \frac{6}{5} = \frac{11}{5}$$

  \item \textbf{Calcul de $C$ :}
  On calcule séparément le numérateur et le dénominateur :
  $$\text{Numérateur : } \frac{2}{3} + \frac{3}{4} = \frac{8}{12} + \frac{9}{12} = \frac{17}{12}$$
  $$\text{Dénominateur : } \frac{5}{6} - \frac{1}{2} = \frac{5}{6} - \frac{3}{6} = \frac{2}{6} = \frac{1}{3}$$
  On effectue ensuite la division des deux fractions :
  $$C = \frac{\frac{17}{12}}{\frac{1}{3}} = \frac{17}{12} \times \frac{3}{1} = \frac{51}{12} = \frac{17}{4}$$
\end{enumerate}

\partie{Correction de la Partie B}
\begin{enumerate}
  \item D'après l'identité remarquable $a^2 - b^2 = (a-b)(a+b)$, on a :
  $$x^2 - 4 = (x - 2)(x + 2)$$

  \item En remplaçant au numérateur par la forme factorisée :
  $$D(x) = \frac{(x - 2)(x + 2)}{x^2} \times \frac{x}{x - 2}$$
  En simplifiant par $(x - 2)$ (valide car $x \neq 2$) et en simplifiant $x$ au numérateur et au dénominateur (valide car $x \neq 0$) :
  $$D(x) = \frac{(x + 2) \times x \times (x - 2)}{x^2 \times (x - 2)} = \frac{x + 2}{x}$$

  \item Pour $x = \frac{1}{3}$ :
  $$D\left(\frac{1}{3}\right) = \frac{\frac{1}{3} + 2}{\frac{1}{3}} = \frac{\frac{1}{3} + \frac{6}{3}}{\frac{1}{3}} = \frac{\frac{7}{3}}{\frac{1}{3}} = \frac{7}{3} \times \frac{3}{1} = 7$$
\end{enumerate}

\end{document}
